\documentclass[11pt,a4paper]{article}

\usepackage[utf8]{inputenc}
\usepackage[T1]{fontenc}
\usepackage{lmodern}
\usepackage[margin=2.5cm]{geometry}
\usepackage{microtype}
\usepackage{titlesec}
\usepackage{parskip}
\usepackage[hidelinks]{hyperref}

\titleformat{\section}{\large\bfseries}{\thesection.}{0.6em}{}
\titlespacing*{\section}{0pt}{1.1em}{0.3em}

\title{\textbf{A Scientific Visualization Platform for Fire Dynamics Simulator Ensembles}}
\author{Internship Progress Report}
\date{}

\begin{document}

\maketitle
\thispagestyle{empty}

\section{Project Overview}

The Fire Dynamics Simulator (FDS) is a software used to simulate how fire, heat, and smoke
move inside buildings. In real studies, researchers usually run many simulations with
different conditions, such as changing the fire size, ventilation, or door openings. The
main challenge is not only running these simulations, but also comparing and understanding
the results.

This project focuses on developing an application that helps researchers visualize and
analyze FDS simulations. It allows them to follow the evolution of physical quantities,
compare different scenarios, and extract useful information from the simulation data.

\section{Initial State of the Project}

At the beginning of the internship, the application already existed but had limited
capabilities. It could display the temperature evolution of a single simulation using a
basic animated heat map. However, it supported only one quantity and one scenario at a time,
with no measurement tools, comparison features, or advanced analysis capabilities.

The objective was to transform it into a modern scientific visualization tool while keeping
the existing structure and improving it step by step instead of rebuilding the whole
application.

\section{Work Performed}

\textbf{Data loading.} Improved the way simulation data is loaded by introducing optimized
reading methods and caching. This made switching between scenarios faster and provided a
better foundation for interactive visualization and analysis.

\textbf{Playback.} Improved the animation system to make it more useful for studying fire
evolution. Users can now move freely through time using a timeline, go frame by frame,
change playback speed, and compare several simulations at the same moment.

\textbf{Scientific visualization.} Added support for multiple physical quantities and new
visual analysis tools. The application now includes side-by-side scenario comparison,
difference views, ensemble views, interactive probes, isotherm and contour visualization,
and extraction of information from volumetric smoke data.

\textbf{Analysis.} Added several tools to study groups of simulations instead of only
individual cases. These include a dataset explorer, statistical summaries, time-series
plots at selected locations, principal component analysis for identifying similar fire
behaviors, forecasting models, and fire-related indicators such as energy release,
smoke-layer height, and safety conditions. The application can also generate reports and
support automated execution through a command-line interface.

\textbf{User interface.} Redesigned the interface to make it easier to use during a
research workflow. The application is now organized around the main tasks: viewing a
simulation, comparing cases, exploring datasets, performing analysis, and exporting
results. Improvements include clearer controls, better visualization of the physical
setup, improved typography, and support for both light and dark modes.

\textbf{Software architecture.} Improved the internal organization of the application by
separating data processing, application logic, and visualization components. A common
system was created for handling different quantities and views, making it easier to add
future features while keeping the software maintainable.

\section{Final Application}

The final result is no longer only a visualization tool. It has become a complete platform
for exploring and analyzing FDS simulation studies.

Researchers can load simulation cases, observe how fire conditions change over time,
compare multiple scenarios, identify important differences, explore datasets, analyze
patterns, and generate reports. The application is also designed to support future FDS
outputs and can be used without always relying on the graphical interface.

\section{Future Work}

The platform provides a strong foundation for future improvements, especially in the area
of artificial intelligence. Some data-driven methods are already included, such as
forecasting and grouping simulations based on their behavior.

Future developments could include more AI-based assistance, such as automatically
identifying fire behaviors, helping researchers interpret simulation results, or exploring
large numbers of simulation parameters more efficiently. These directions will be studied
and developed in future work.

\section{Conclusion}

This internship transformed a simple FDS visualization application into a more complete
scientific analysis platform. The project introduced new visualization, comparison,
measurement, and analysis capabilities while keeping the original system and improving it
progressively.

The final application provides researchers with a more efficient way to study fire
simulations and creates a solid foundation for future developments, including the
integration of more advanced artificial intelligence techniques.

\end{document}